\documentclass[10pt,a4paper]{amsart}
\usepackage[margin=19mm]{geometry}
\usepackage{amsmath,amssymb,mathtools}
\usepackage[scr=rsfs,cal=euler]{mathalfa}
\usepackage{xfrac}
\usepackage[unicode,hidelinks,bookmarks=false]{hyperref}
\newcommand{\ie}{i.e.\ }
\newcommand{\cf}{cf.\ }
\newcommand{\resp}{respectively\ }
\newcommand{\Subsection}{\subsection}
\providecommand{\nobreakdash}{\nobreak\hbox{-}}
\setlength{\emergencystretch}{2em}
\multlinegap0pt
\usepackage{algorithm}\usepackage{algpseudocode}
\usepackage{amssymb}
\newtheorem{lemma}{Lemma}
\newcommand{\E}{\mathbb{E}} % Expectation
\newcommand{\N}{\mathbb{N}} % natural numbers
\newcommand{\dirE}{\vec{E}}
\DeclareMathOperator{\dom}{dom}
\DeclareMathOperator{\outdeg}{outdeg}
\DeclareMathOperator{\var}{var}
\title{Entropy and the growth rate of universal covering trees}
\author{Idan Eisner, Shlomo Hoory}
\date{Source: arXiv:2410.10337v3 (CC BY 4.0)}
\begin{document}
\maketitle

\noindent\emph{Source and licence.} Entropy and the growth rate of universal covering trees. Version arXiv:2410.10337v3 (CC BY 4.0), distributed by arXiv under the Creative Commons Attribution 4.0 International licence, http://creativecommons.org/licenses/by/4.0/, as recorded in the arXiv licence metadata for this exact version and shown on the abstract page https://arxiv.org/abs/2410.10337v3. Full licence notice: \texttt{licenses/arxiv-2410.10337-CC-BY-4.0.txt}.\par\smallskip

\noindent\textit{Editorial scope.} Counted unit: the article's lemma lemma:variance\_lb (source0.tex lines 756--758). The article's complete original proof of it is delivered with it (source0.tex lines 765--825, one \texttt{proof} environment of 61 lines, which the article places away from the statement). No mathematics is written, repaired or re-derived: the statement and the proof are the article's own lines, in the article's own notation, and no retained line is substituted or re-flowed.  Retained uncounted as the source-local context the proof needs: lines 720--726.  Nothing was omitted from any retained span, so the package carries no omission record.  No retained line contains a citation, so the wrapper carries no bibliography and no bibliography entry.  No retained line contains a figure environment, an image-inclusion command or drawing code, so no image is delivered and the package declares no asset dependency.  Editorial formatting only, in addition: an AMS \texttt{amsart} preamble that restates the article's own declarations and macros used by the retained lines, namely the environments \texttt{lemma} on separate counters, together with the standard packages those lines need.  The retained environments print in this document as Lemma 1 in order of appearance, whereas the article prints the same items with the numbering of its own full document.  The complete native source of the article as distributed in the arXiv e-print for this exact version is bundled unchanged as \texttt{sources/166785/source0.tex}.
\begin{lemma}\label{lemma:two_cycles_variance}
    Given an NB-irreducible graph $G$ with $\rho > \Lambda$, there exist two non-backtracking cycles $C_1$ and $C_2$ with $|C_1| = |C_2|$, 
    a common edge in $C_1 \cap C_2$, but different average out-degree:
    \begin{equation*}
        \left[ \prod_{e \in C_1} \outdeg(e) \right]^{\dfrac{1}{|C_1|}} < \left[ \prod_{e \in C_2} \outdeg(e) \right]^{\dfrac{1}{|C_2|}}.
    \end{equation*}
\end{lemma}

\begin{lemma}\label{lemma:variance_lb}
    If $G$ is an NB-irreducible graph with $\rho > \Lambda$ then there is a constant $c=c(G)$ so that $\var[R_\ell] \geq c \cdot \ell$ for all $\ell$.
\end{lemma}

\begin{proof}[Proof of Lemma~\ref{lemma:variance_lb}]
    Given an NB-irreducible graph $G$ with $\rho > \Lambda$,
    by Lemma~\ref{lemma:two_cycles_variance} 
    there are two cycles $C_1$ and $C_2$ of the same length $\ell_0=|C_1| = |C_2|$, 
    with different average out-degree, and with some edge $f$ in their intersection.
    Recall that:
    \begin{equation*}
        \var[R_\ell] = \E_{\omega \in \Omega_\ell}\left[ {(R_\ell(\omega) - \ell \log_2\Lambda)}^2\right].
    \end{equation*}
    In order to lower-bound the variance, we use the exposure method. 
    We expose some of the edges of $\omega = (e_0,e_1,\ldots,e_\ell)$
    and argue that the conditional expectation of ${(R_\ell(\omega) - \ell \log_2\Lambda)}^2$ given the exposed edges is $\Omega(\ell)$ with high probability.
    Define the mapping $\phi:\dirE\to\N$, where $\phi(e)$ is the length of the shortest NB-path from $e$ to $f$ if $e \neq f$ and $\phi(f) = \ell_0$.
    Also, let $\ell_{\max}$ be the maximum of $\phi(e)$ on $e \in \dirE$.
    We build the partial exposure mapping $\chi$ from $\{0,\ldots,\ell\}$ to $\dirE$ by the following randomized procedure:
    
    \begin{center}
        \fbox{\begin{minipage}{7 cm}
        \begin{algorithmic}
            \State $\chi \gets $ the empty mapping
            \State $i \gets 0$
            \While{$i < \ell$}
                \State Expose $e_i$ and set $\chi(i) \gets e_i$
                \State $i \gets i + \phi(e_i)$
            \EndWhile
            \State Expose $e_\ell$ and set $\chi(\ell) \gets e_\ell$
        \end{algorithmic}
        \end{minipage}}        
    \end{center}
    
    The domain of $\chi$ is the set of exposed edge indices, $i_1=0 < i_2 < \cdots < i_k=\ell$.
    Therefore, conditioned on $\chi$, one may write $R_\ell$ as the sum of independent random variables:
    \begin{equation*}
        R_\ell = \sum_{i=0}^{\ell-1} \log_2 \outdeg(e_i) = \sum_{j=1}^{k-1} R(i_{j+1}-i_j,e_{i_j},e_{i_{j+1}}),
    \end{equation*}
    where $R(\ell',e',e'')$ is the random variable counting the number of random bits used by a length $\ell'$ NBRW, conditioned on it beginning at $e'$ and ending at $e''$.
    Therefore:
    \begin{equation*}
        \var[R_\ell | \chi]
           = \sum_{j=1}^{k-1} \var\!\left[R(i_{j+1}-i_j,e_{i_j},e_{i_{j+1}})\right]
           \geq |I| \cdot \var\!\left[R(\ell_0,f,f)\right],
    \end{equation*}
    where $I = I(\chi) = \{i \in  \dom(\chi) :\, i+\ell_0 \in \dom(\chi) \mbox{ and } \chi(i) = \chi(i+\ell_0) = f\}$.
    As the number of exposed edges is lower bounded by $\ell/\ell_{\max} = \Omega(\ell)$,
    it follows that $|I|$ is lower bounded by a binomial random variable $X \sim B(\ell/(2\ell_{\max}),p_{\min})$, 
    where $p_{\min}$ is a positive constant defined as the minimum of $(\Pi^{\phi(e)})_{e,f} \cdot (\Pi^{\ell_0})_{f,f}$ on $e \in \dirE$. 
    Therefore, $\E[|I|] \geq \E[X] = c_1 \ell$, where $c_1 = p_{\min}/(2 \ell_{\max})$.
    The Chernoff bound on the lower tail of $X$ yields $\Pr[|I| \leq c_1 \ell/2] \leq \Pr[X \leq c_1 \ell/2] \leq \exp[-c_1 \ell/2]$, which is $o(1)$ as $\ell$ grows.

    Recall that by our assumption on the cycles $C_1, C_2$, the variance of $R(\ell_0,f,f)$ is some positive constant $c_2$.
    Therefore:
    \begin{eqnarray*}
        \var[R_\ell]
        &=& \E\!\left[{(R_\ell(\omega) - \ell \log_2\Lambda)}^2\right]
        = \sum_\chi \E\!\left[{(R_\ell(\omega) - \ell \log_2\Lambda)}^2 \Big\vert \chi\right] \cdot \Pr[\chi] \\
        &\geq& \sum_\chi |I(\chi)| \cdot \var\!\left[R(\ell_0,f,f)\right] \cdot \Pr[\chi] \\
        &\geq& \frac{c_1 \ell}{2} \cdot \var\!\left[R(\ell_0,f,f)\right] \cdot \Pr\!\left[|I| \geq \frac{c_1 \ell}{2}\right] \\
        &\geq& \frac{c_1 c_2 \ell}{2} \cdot \Pr\!\left[|I| \geq \frac{c_1 \ell}{2}\right]
        = \Omega(\ell).
    \end{eqnarray*}
\end{proof}

\end{document}
